% 二元选择 max x1+x2, 2x1+2x2<=3, x in {0,1}^2。(a) 线性规划松弛 {0<=x<=1, x1+x2<=3/2}（蓝）包含精确凸包 conv{(0,0),(1,0),(0,1)}（绿）；
% 松弛最优面为(1,1/2)-(1/2,1)，值3/2；整数最优值1。(b) 从(3/4,3/4)四舍五入到(1,1)不可行，向下取整到(0,0)值0；从(1,1/2)修复到(1,0)值1。
\begin{tikzpicture}[x=30mm,y=30mm,
  note/.style={font=\FigureNoteFont\fontsize{7.5}{9}\selectfont,text=FigureMuted},
  ttl/.style={font=\fontsize{8}{10}\selectfont,anchor=north},
  arr/.style={-{Stealth[length=2mm]},line width=.8pt}]
  \begin{scope}
    \fill[FigureBlueFill] (0,0)--(1,0)--(1,0.5)--(0.5,1)--(0,1)--cycle;
    \fill[FigureGreenFill] (0,0)--(1,0)--(0,1)--cycle;
    \draw[FigureBlue,line width=.9pt] (0,0)--(1,0)--(1,0.5)--(0.5,1)--(0,1)--cycle;
    \draw[FigureGreen,line width=.9pt] (1,0)--(0,1);
    \draw[FigureRed,line width=2pt] (1,0.5)--(0.5,1);
    \draw[->,FigureStroke,line width=.6pt] (-0.06,0)--(1.3,0) node[right]{$x_1$};
    \draw[->,FigureStroke,line width=.6pt] (0,-0.06)--(0,1.25) node[above]{$x_2$};
    \foreach \p in {(0,0),(1,0),(0,1)}{\fill[FigureBlue] \p circle (2pt);}
    \draw[FigureRed,line width=.9pt] (1,1) circle (2pt);
    \fill[FigureRed] (0.75,0.75) ++(-1.6pt,-1.6pt) rectangle ++(3.2pt,3.2pt);
    \node[note,anchor=south west] at (1.02,1.0) {$(1,1)$不可行};
    \node[note,anchor=south west,text=FigureRed!80!FigureInk] at (0.78,0.68) {值$3/2$};
    \node[note,text=FigureGreen!75!FigureInk] at (0.3,0.28) {精确凸包};
    \node[note,text=FigureBlue!80!FigureInk,anchor=west] at (0.55,0.38) {松弛};
    \node[ttl] at (0.6,-0.12) {(a) 外逼近松弛};
  \end{scope}
  \begin{scope}[xshift=50mm]
    \fill[FigureBlueFill] (0,0)--(1,0)--(1,0.5)--(0.5,1)--(0,1)--cycle;
    \draw[FigureBlue!60,line width=.6pt] (0,0)--(1,0)--(1,0.5)--(0.5,1)--(0,1)--cycle;
    \draw[->,FigureStroke,line width=.6pt] (-0.06,0)--(1.3,0) node[right]{$x_1$};
    \draw[->,FigureStroke,line width=.6pt] (0,-0.06)--(0,1.25) node[above]{$x_2$};
    \foreach \p in {(0,0),(1,0),(0,1)}{\fill[FigureBlue] \p circle (2pt);}
    \draw[FigureRed,line width=.9pt] (1,1) circle (2pt);
    \fill[FigureRed] (0.75,0.75) ++(-1.6pt,-1.6pt) rectangle ++(3.2pt,3.2pt);
    \fill[FigureRed] (1,0.5) ++(-1.6pt,-1.6pt) rectangle ++(3.2pt,3.2pt);
    \draw[arr,FigureRed] (0.78,0.78)--(0.96,0.96);
    \draw[arr,FigureMuted] (0.72,0.72)--(0.04,0.04);
    \draw[arr,FigureGreen] (1.03,0.47)--(1.03,0.05);
    \node[note,anchor=east,text=FigureRed!80!FigureInk] at (0.95,1.08) {四舍五入};
    \node[note,anchor=north west] at (0.3,0.33) {向下取整，值$0$};
    \node[note,anchor=west,text=FigureGreen!75!FigureInk] at (1.06,0.26) {修复，值$1$};
    \node[ttl] at (0.6,-0.12) {(b) 恢复原候选};
  \end{scope}
\end{tikzpicture}
